\documentclass[10pt]{amsart}
\usepackage[T1]{fontenc}
\usepackage{lmodern}
\usepackage{amsmath,amssymb,amsthm}
\usepackage{xcolor}
\usepackage[colorlinks=true,linkcolor=blue!55!black,citecolor=blue!55!black,urlcolor=blue!55!black]{hyperref}
\newtheorem{theorem}{Theorem}
\newtheorem{proposition}[theorem]{Proposition}
\title{A Counterexample to the Positivity Conjecture for OEIS A071532}
\author{Xiaojian Zeng}
\address{PBC School of Finance, Tsinghua University, 100190}
\email{xiaojianzeng@foxmail.com}
\date{September 2, 2026}
\hypersetup{pdftitle={A Counterexample to the Positivity Conjecture for OEIS A071532},pdfauthor={Xiaojian Zeng}}
\begin{document}
\begin{abstract}
Let $S_n=-\sum_{k=1}^n(-1)^{\lfloor(3/2)^k\rfloor}$.
The recorded conjecture that $S_n>0$ for every $n>2$ is false. Exact
integer enumeration gives $S_{331522}=0$; moreover, $331522$ is the first
violating index. We give the elementary binary-digit reduction and a small,
independently rerunnable verifier that checks every prefix through $10^6$.
No floating-point arithmetic or unverified solver output is used.
\end{abstract}
\maketitle
\section{Statement and provenance}
OEIS A071532 records
\begin{equation}\label{eq:def}
 S_n=-\sum_{k=1}^n(-1)^{\lfloor(3/2)^k\rfloor}
\end{equation}
and asks whether $S_n>0$ \cite{oeis71532}. Since its displayed values include
$S_2=0$, the precise version is the assertion for every integer $n>2$.
The pinned Formal Conjectures source states exactly this endpoint, but leaves
the theorem unproved \cite{formal71532}. Our searches through the cutoff date
found no earlier published resolution; this is a bounded literature statement,
not a claim that every possible source has been exhausted.

One comment in the current OEIS record has a harmless indexing error. If
$c(n)$ counts odd quotients over $0\leq k\leq n$, then the $k=0$ quotient is
odd, and hence $S_n=2c(n)-n-2$, rather than $2c(n)-n$. The defining formula
\eqref{eq:def} and the conjecture are unaffected.

\section{Exact reduction and counterexample}
Put $q_k=\lfloor3^k/2^k\rfloor$. Euclidean division gives
\[
 3^k=2^kq_k+r_k,\qquad 0\leq r_k<2^k.
\]
Consequently $q_k\bmod2$ is precisely bit $k$ (numbering the least
significant bit as zero) of the integer $3^k$. Thus $S_0=0$ and the exact
recurrence is
\begin{equation}\label{eq:update}
 S_k=S_{k-1}+\begin{cases}1,&((3^k\mathbin{\mathrm{div}}2^k)\bmod2)=1,\\
 -1,&((3^k\mathbin{\mathrm{div}}2^k)\bmod2)=0.
 \end{cases}
\end{equation}
\begin{theorem}\label{thm:counterexample}
The assertion $S_n>0$ for every integer $n>2$ is false. In fact,
\[
 S_{331522}=0,
\]
and $S_n>0$ for $3\leq n<331522$.
\end{theorem}
\begin{proof}
Starting with $p=1$ and $S=0$, for each $k=1,\ldots,331522$ replace
$p$ by $3p$, extract the integer bit $(p\mathbin{\mathrm{div}}2^k)\bmod2$,
and update $S$ by \eqref{eq:update}. This is a complete enumeration using
only integer multiplication, division by powers of two, parity, and addition.
The supplied verifier recomputes every update and reports the first
nonpositive target prefix as $331522$, with value zero. A second evaluator
checks the first $10000$ bits from $3^k\bmod2^{k+1}$ and from a fresh direct
quotient. Since all earlier target indices are visited and tested, the first
claim follows; the exact equality at $331522$ disproves strict positivity.
\end{proof}

For local inspectability, the last eleven updates around the first failure are
\begingroup\scriptsize\setlength{\arraycolsep}{2.2pt}
\[
\begin{array}{c|rrrrrrrrrrr}
k&331515&331516&331517&331518&331519&331520&331521&331522&331523&331524&331525\\\hline
S_k&7&6&5&4&3&2&1&0&-1&-2&-3.
\end{array}
\]
\endgroup
In each of these eleven steps the extracted bit is zero. This short table is
illustrative; the rerun, rather than the table, certifies all preceding steps.

\begin{proposition}
The same exact run certifies all prefixes through $10^6$. It gives
$S_{10^6}=580$ and a minimum value $-359$, attained at $n=371843$, over
$3\leq n\leq10^6$.
\end{proposition}

\section{Reproducibility and limitations}
From the source root the certificate is rerun by
\begin{quote}\small
\texttt{python3 certificate/certifier\_verify.py}\\
\texttt{\phantom{python3 }certificate/certifier\_claim.json}
\end{quote}
The recorded run used CPython 3.10.12 and only its standard library. There is
no randomness, floating-point arithmetic, SAT/SMT solver, or proof assistant.
The code rejects malformed claims and independently recomputes every update
and all stated statistics. The finite behavior after the first
failure is not needed for Theorem~\ref{thm:counterexample}; it is reported only
as an additional reproducibility check. The computation does not address the
separate square-root or log-squared asymptotic questions on the OEIS page.

The exact enumeration, certificate generation, source build, and document
checks were performed with standard programmatic tools. All mathematical
claims needed for the disproof reduce to the displayed integer recurrence and
can be rerun from the archived source and data.
\bibliographystyle{amsplain}
\begingroup\sloppy
\bibliography{references}
\endgroup
\end{document}
